\documentclass[french,11pt]{book}
\usepackage{teach}
\usepackage[french]{babel}
\usepackage{amsmath,amssymb}
\usepackage{array}
\newcommand{\R}{\mathbb{R}}
\begin{document}
\begin{center}
{\Large \textbf{Exercices -- Loi binomiale}}\\[2mm]
Terminale technologique
\end{center}

\section*{Exercices}
\begin{enumerate}
  \item On répète 8 fois une expérience de Bernoulli de probabilité de succès $p=0,3$. Définir la variable aléatoire $X$ associée et préciser sa loi.
  \item Pour $X\sim\mathcal{B}(10;0,4)$, calculer à la calculatrice $P(X=3)$, $P(X\leq 2)$ et $P(X\geq 5)$.
  \item Dans une production, 5\% des pièces sont défectueuses. On prélève 20 pièces de manière assimilée à des tirages indépendants. Déterminer la probabilité d'obtenir exactement deux pièces défectueuses.
  \item Pour $X\sim\mathcal{B}(n;p)$, rappeler les formules de l'espérance et de l'écart type, puis les appliquer à $n=50$ et $p=0,12$.
\end{enumerate}

\end{document}
